\exercisesection{}

In Exercises 1 to 5 exclusively, K denotes a complete non-Archimedean non-discrete valued field, whose valuation is denoted by \(x \mapsto |x|\) . We also denote by \(\mathrm{G} \subset \mathbf{R}_{+}^{*}\) the image of \(\mathrm{K}^{*}\) under the map \(x \mapsto |x|\) . The Banach spaces considered are Banach spaces over \(\mathrm{K}\) .

\begin{enumerate}
\def\labelenumi{\arabic{enumi})}
\tightlist
\item
  \begin{enumerate}
  \def\labelenumii{\alph{enumii})}
  \tightlist
  \item
    Let E be a Banach space over K. There exists a norm \(x \mapsto \|x\|\) on E, defining the topology of E, such that \(\|x\| \in \overline{G}\) for all \(x \in E\) . Such a norm will be called round. If the valuation of K is not discrete, then every norm defining the topology of E is round.
  \end{enumerate}
\end{enumerate}

\begin{enumerate}
\def\labelenumi{\alph{enumi})}
\setcounter{enumi}{1}
\item
  Let E and F be Banach spaces over K. We equip the space \(\mathcal{L}(\mathrm{E};\mathrm{F})\) with the norm \(\| u\| = \sup_{x\in \mathrm{E}- \{0\}}\| u(x)\| /\| x\|\) . The space \(\mathcal{L}(\mathrm{E};\mathrm{F})\) is a Banach space over K. If the norm on E is round, then \(\| u\| = \sup_{\| x\| \leqslant 1}\| u(x)\|\) . If the norm on F is also round, then the norm on \(\mathcal{L}(\mathrm{E};\mathrm{F})\) is round. We denote \(\mathrm{E}' = \mathcal{L}(\mathrm{E};\mathrm{K})\) and we say that \(\mathrm{E}'\) is the dual space of E.
\item
  Let I be a set, F a Banach space over K, and \(\mathcal{C}_{0}(\mathrm{I},\mathrm{F})\) the K-vector space of functions \(f\colon I\to F\) such that \(\|f(i)\|\) tends to 0 along the filter of complements of finite subsets of I. Equipped with the norm \(\|f\|=\sup_{i\in I}\|f(i)\|\) , the space \(\mathcal{C}_{0}(\mathrm{I},\mathrm{F})\) is a Banach space over K, whose norm is round if that of F is. We denote \(\mathcal{C}_{0}(\mathrm{I})=\mathcal{C}_{0}(\mathrm{I},\mathrm{K})\) .
\end{enumerate}

\(d\) ) Let I be a set. For every \(i\in \mathrm{I}\) , we denote by \(\varphi_{i}\in \mathcal{C}_{0}(\mathrm{I})\) the characteristic function of \(\{i\}\) . Let F be a Banach space over K. The map \(u \mapsto (u(\varphi_i))_{i \in \mathrm{I}}\) is an isomorphism from the Banach space \(\mathcal{L}(\mathcal{C}_0(\mathrm{I});\mathrm{F})\) to the Banach space \(\mathcal{B}(\mathrm{I};\mathrm{K})\) (EVT, I, p. 4, example; cf. EVT, I, p. 23, exercise 7).

\begin{enumerate}
\def\labelenumi{\alph{enumi})}
\setcounter{enumi}{4}
\item
  Assume that the valuation of K is discrete. Let E be a Banach space over K whose norm is round. There exists a set I such that E is isometrically isomorphic to the Banach space \(\mathcal{C}_{0}(\mathrm{I})\) (cf. EVT, I, p. 23, exercise 7).
\item
  Assume that the valuation of K is discrete. Let E be a Banach space over K equipped with a round norm and F a closed vector subspace of E. There exists a continuous projector of E whose image is F and which has norm \(\leqslant 1\) .
\end{enumerate}

\begin{enumerate}
\def\labelenumi{\arabic{enumi})}
\setcounter{enumi}{1}
\tightlist
\item
  Let E and F be Banach spaces over K. We say that \(u \in \mathcal{L}(\mathrm{E};\mathrm{F})\) is completely continuous if it belongs to the closure in \(\mathcal{L}(\mathrm{E};\mathrm{F})\) of the space \(\mathcal{L}^{\mathrm{f}}(\mathrm{E};\mathrm{F})\) of finite-rank linear maps. We denote by \(\mathcal{L}^{\mathrm{c}}(\mathrm{E};\mathrm{F})\) the space of completely continuous linear maps from E to F, and we set \(\mathcal{L}^{\mathrm{c}}(\mathrm{E}) = \mathcal{L}^{\mathrm{c}}(\mathrm{E};\mathrm{E})\) . a) The space \(\mathcal{L}^{\mathrm{c}}(\mathrm{E})\) is a closed two-sided ideal of \(\mathcal{L}(\mathrm{E})\) . Let \(u \in \mathcal{L}(\mathrm{E};\mathrm{F})\) . The map \(\ell \mapsto \ell \circ u\) is a continuous linear map from \(F'\) to \(E'\) , denoted \({}^{t}u\) . If \(u \in \mathcal{L}^{\mathrm{c}}(\mathrm{E};\mathrm{F})\) , then \({}^{t}u \in \mathcal{L}^{\mathrm{c}}(\mathrm{F}'; \mathrm{E}')\) .
\end{enumerate}

\begin{enumerate}
\def\labelenumi{\alph{enumi})}
\setcounter{enumi}{1}
\item
  Let J be a set and \(\mathrm{F} = \mathcal{C}_{0}(\mathrm{J})\) . Let \(u \in \mathcal{L}(\mathrm{E};\mathrm{F})\) . For every \(j \in J\) , the map \(\ell_{j} \colon x \mapsto u(x)(j)\) is an element of \(E'\) . We have \(\|u\| = \sup_{j \in J} \| \ell_{j} \|\) and the map \(u \mapsto (\ell_{j})_{j \in J}\) induces, by passage to subspaces, an isomorphism from \(\mathcal{L}^{\mathrm{c}}(\mathrm{E};\mathrm{F})\) to the Banach space \(\mathcal{C}_{0}(\mathrm{J};\mathrm{E}')\) .
\item
  If K is locally compact, then \(u \in \mathcal{L}(\mathrm{E}; \mathrm{F})\) is completely continuous if and only if the image under u of every bounded subset of E is relatively compact in F. (One may compare this result with prop. 11 of III, p. 16.)
\end{enumerate}

\begin{enumerate}
\def\labelenumi{\arabic{enumi})}
\setcounter{enumi}{2}
\tightlist
\item
  We denote by A the subring of K consisting of elements \(x \in \mathrm{K}\) such that \(|x| \leqslant 1\) ; it is a local ring whose maximal ideal consists of the \(x \in \mathrm{A}\) such that \(|x| < 1\) .
\end{enumerate}

Let I be a set and E the Banach space \(\mathcal{C}_{0}(\mathrm{I})\) (cf. question c) of exercise 1). For \(i \in I\) , denote by \(\varphi_{i} \in E\) the characteristic function of \(\{i\}\) . Let \(u \in \mathcal{L}^{\mathrm{c}}(\mathrm{E})\) such that \(\|u\| \leqslant 1\) .

\begin{enumerate}
\def\labelenumi{\alph{enumi})}
\item
  The subspace \(E_{0}\) of elements \(x \in E\) such that \(\|x\| \leqslant 1\) is stable under u.
\item
  For every nonzero ideal a of A, the linear map u defines, by passage to quotients, an endomorphism \(u_{a}\) of the A/a-module \(E_{a} = E_{0}/aE_{0}\) . This A/a-module is free and the image of \(u_{a}\) is contained in a finitely generated submodule.
\item
  There exists a unique formal series \(f_{u} \in A[[t]]\) such that, for every nonzero ideal a of A, the image of \(f_{u}\) in \((A/\mathfrak{a})[[t]]\) is equal to the polynomial \(\det(1 - tu_{\mathfrak{a}})\) defined in A, VIII, p. 455.
\end{enumerate}

\(d)\) Let \(v \in \mathcal{L}^{\mathrm{c}}(\mathrm{E})\) and let \(c \in \mathrm{K}^{*}\) such that \(\|cv\| \leqslant 1\) . The formal series \(f_{cv}\) is independent of the choice of \(c\) satisfying this property.

We shall denote by \(\det(1-tv)\) the formal series \(f_v\) for all \(c \in \mathbf{K}^*\) such that \(\|cv\| \leqslant 1\) .

¶ 4) We retain the notation from the previous exercise.

\begin{enumerate}
\def\labelenumi{\alph{enumi})}
\item
  Let \(v \in \mathcal{L}^c(\mathrm{E})\) . The radius of convergence of the formal series \(\det(1 - tv)\) is infinite.
\item
  Let \((v_{n})_{n\in\mathbf{N}}\) be a sequence in \(\mathcal{L}^{c}(\mathrm{E})\) converging to \(v\in\mathcal{L}^{c}(\mathrm{E})\) . The sequence of formal series \(\det(1-tv_{n})\) converges to \(\det(1-tv)\) for the topology of simple convergence of the coefficients of formal series.
\item
  Let \(v \in \mathcal{L}^{c}(E)\) . For every extension L of K and every \(x \in L\) , we denote by \(\det(1 - xv)\) the value at x of the formal power series \(\det(1 - tv)\) , which is well defined by a). We have \(\det(1 + u)\det(1 + v) = \det(1 + u + v + u \circ v)\) .
\item
  Let J be a set and \(\mathrm{F} = \mathcal{C}_{0}(\mathrm{J})\) . For all \(u \in \mathcal{L}^{\mathrm{c}}(\mathrm{E};\mathrm{F})\) and \(v \in \mathcal{L}^{\mathrm{c}}(\mathrm{F};\mathrm{E})\) , we have \(\det(1 - tu \circ v) = \det(1 - tv \circ u)\) as formal power series.
\item
  Let \(v \in \mathcal{L}^{\mathrm{c}}(\mathrm{E})\) . We denote by \(\operatorname{Tr}(v)\) the negative of the coefficient of t in \(\det(1 - tu)\) . If K is of characteristic zero, then
\end{enumerate}

\[
\det (1 - t u) = \exp \left(- \sum_ {m \geqslant 1} \frac {\operatorname{Tr} \left(u ^ {m}\right)}{m} t ^ {m}\right)
\]

as formal power series.

¶ 5) We keep the notations of the preceding exercise. Let \(v \in \mathcal{L}^{\mathrm{c}}(\mathrm{E})\) .

\begin{enumerate}
\def\labelenumi{\alph{enumi})}
\item
  The formal power series \(\det(1 - tv)\sum_{k\geqslant 0}t^{k}v^{k}\) in \(\mathcal{L}(\mathrm{E})[[t]]\) has infinite radius of convergence.
\item
  Let \(\lambda \in \mathrm{K}\) . The endomorphism \(1_{\mathrm{E}} - \lambda v\) of E is invertible if and only if \(\det(1 - \lambda v) \neq 0\) .
\item
  Let \(\lambda \in K\) such that \(\det(1 - \lambda v) = 0\) . There exist unique closed subspaces \(I_{\lambda}\) and \(N_{\lambda}\) of E, stable under \(v\) such that
\end{enumerate}

\begin{enumerate}
\def\labelenumi{(\roman{enumi})}
\item
  The space E is the topological direct sum of \(I_{\lambda}\) and \(N_{\lambda}\) ;
\item
  The endomorphism of \(I_{\lambda}\) deduced from \(1_{E} - \lambda v\) by passing to subspaces is invertible;
\item
  The endomorphism of \(N_{\lambda}\) deduced from \(1_{E} - \lambda v\) by passing to subspaces is nilpotent.
\end{enumerate}

\(d)\) Let \(h \geqslant 1\) be the order of the zero \(\lambda\) of det(1 - tv). We have \(\dim(N_{\lambda}) = h\) .

The result of this exercise is to be compared with those of No. 5 of III, p. 77, which concern Banach spaces over R or C.

In what follows, K denotes the field R or C, and all topological vector spaces are over K.

\begin{enumerate}
\def\labelenumi{\arabic{enumi})}
\setcounter{enumi}{5}
\tightlist
\item
  Let E be the space \(C^{N}\) endowed with the topology associated with the sequence of semi-norms \((p_{n})_{n\in\mathbf{N}}\) given by \(p_{n}(x)=|x_{n}|\) for every element \(x=(x_{n})_{n\in\mathbf{N}}\) of E.
\end{enumerate}

\begin{enumerate}
\def\labelenumi{\alph{enumi})}
\item
  The space E is a Fréchet space.
\item
  The set \(\mathcal{L}^{\mathrm{c}}(\mathrm{E})\) of compact endomorphisms of E is not closed in \(\mathcal{L}(\mathrm{E})\) .
\end{enumerate}

\begin{enumerate}
\def\labelenumi{\arabic{enumi})}
\setcounter{enumi}{6}
\item
  Let E and F be locally convex topological vector spaces. If F is metrizable, then the image of every compact linear map \(E \rightarrow F\) is of countable type.
\item
  Construct an example of a non-compact endomorphism \(u\) of a Hilbert space E such that \(u^2\) is compact.
\item
  For every pair \((x,y)\in]0,1]^{2}\) , there exists an integer \(n\geqslant1\) and an integer k such that \(0\leqslant k\leqslant2^{n-1}-1\) so that
\end{enumerate}

\[
2 ^ {- n} <   x \leqslant 2 ^ {- n + 1}, \quad \frac {2 k}{2 ^ {n}} <   y \leqslant \frac {2 k + 2}{2 ^ {n}}.
\]

We then define N: ]0,1]×]0,1] → C by setting

\[
\mathrm{N} (x, y) = \left\{ \begin{array}{l l} 2 & \text { if } \frac {2 k}{2 ^ {n}} <   y \leqslant \frac {2 k + 1}{2 ^ {n}} \\ 0 & \text { if } \frac {2 k + 1}{2 ^ {n}} <   y \leqslant \frac {2 k + 2}{2 ^ {n}}. \end{array} \right.
\]

\begin{enumerate}
\def\labelenumi{\alph{enumi})}
\tightlist
\item
  Show that the formula
\end{enumerate}

\[
u (f) (x) = \int_ {[ 0, 1 ]} \mathrm{N} (x, y) f (y) d y
\]

defines a continuous linear map u on \(L^{1}([0,1])\) with respect to Lebesgue measure.

\begin{enumerate}
\def\labelenumi{\alph{enumi})}
\setcounter{enumi}{1}
\item
  Show that the endomorphism \(u\) is not compact.
\item
  The endomorphism \(u^2\) is compact.
\end{enumerate}

\(d)\) Let \((x_{n})_{n\in \mathbf{N}}\) be a bounded sequence in \(\mathrm{L}^1 ([0,1])\) . The sequence \((u(x_n))_{n\in \mathbf{N}}\) admits a weakly convergent subsequence.

\begin{enumerate}
\def\labelenumi{\alph{enumi})}
\setcounter{enumi}{4}
\item
  Let \((x_{n})_{n\in \mathbf{N}}\) be a weakly convergent sequence in \(\mathrm{L}^1 ([0,1])\) . The sequence \((u(x_{n}))_{n\in \mathbf{N}}\) converges in \(\mathrm{L}^1 ([0,1])\) .
\item
  Let E be the topological vector space whose elements are the functions \(f\colon[0,1]\to\mathbf{C}\) which are bounded and measurable, equipped with the finest topology that is stronger than the topology of \(\mathrm{L}^{2}([0,1])\) and the weak topology relative to \(\mathrm{L}^{1}([0,1])\) . Then the linear map v defined by
\end{enumerate}

\[
v (f) (x) = \int_ {[ 0, 1 ]} \overline {{\mathrm{N} (y , x)}} f (y) d y
\]

is a compact endomorphism of E. Its transpose \(^{t}v\) is not compact for the strong topology on E'.

\begin{enumerate}
\def\labelenumi{\arabic{enumi})}
\setcounter{enumi}{9}
\tightlist
\item
  Let p and r be real numbers such that \(1 \leqslant p < r\) . Let u be a continuous linear map from \(\ell_{\mathrm{K}}^{r}(\mathbf{N})\) into \(\ell_{\mathrm{K}}^{p}(\mathbf{N})\) .
\end{enumerate}

We denote by \(r_{n}\) (resp. \(p_{n}\) ) the projection

\[
(x _ {i}) _ {i \in \mathbf {N}} \mapsto (x _ {0}, \ldots , x _ {n - 1}, 0, 0, \ldots)
\]

onto \(\ell_{\mathrm{K}}^{r}(\mathbf{N})\) (resp. on \(\ell_{\mathrm{K}}^{p}(\mathbf{N})\) ). We denote by \(\| \cdot \| _p\) (resp. \(\| \cdot \| _r\) ) the norm on \(\ell_{\mathrm{K}}^{p}(\mathbf{N})\) (resp. on \(\ell_{\mathrm{K}}^{r}(\mathbf{N})\) ).

\begin{enumerate}
\def\labelenumi{\alph{enumi})}
\item
  Show that for every \(n \in \mathbf{N}\) , we have \(\|p_n u(1 - r_m)\| \to 0\) as \(m \to +\infty\) where 1 denotes the identity morphism of \(\ell_{\mathrm{K}}^{r}(\mathbf{N})\) .
\item
  For \(n\) , \(m \in \mathbf{N}\) , let \(u_{n,m} = (1 - p_n)u(1 - r_m)\) . Let \(\alpha = (\alpha_n)_{n \in \mathbf{N}}\) be a sequence of strictly positive real numbers such that \(\|\alpha\|_r = 1\) and \(\alpha\) does not belong to \(\ell_{\mathrm{K}}^p(\mathbf{N})\) and let \((\varepsilon_n)\) be a sequence of strictly positive real numbers such that the series \(\sum \varepsilon_n\alpha_n\) converges.
\end{enumerate}

If \(u\) is not compact, show that there exist \(\delta > 0\) , strictly increasing sequences of integers \((n_i)_{i \in \mathbf{N}}\) and \((m_i)_{i \in \mathbf{N}}\) , and two sequences \((x_i)_{i \in \mathbf{N}}\) of elements of \(\ell_{\mathrm{K}}^r(\mathbf{N})\) and \((y_i)_{i \in \mathbf{N}}\) of elements of \(\ell_{\mathrm{K}}^p(\mathbf{N})\) with the following properties:

\begin{enumerate}
\def\labelenumi{(\roman{enumi})}
\item
  One has \(\|x_{i}\|_{r}=1\) for all \(i\in N\) ;
\item
  The sets \(S_{i} = \{n \in N \mid x_{i,n} \neq 0\}\) are pairwise disjoint and the sets \(T_{i} = \{n \in N \mid y_{i,n} \neq 0\}\) are pairwise disjoint;
\item
  One has \(\| y_i\| _p\geqslant \delta\)
\item
  One has \(\|u(x_{i})\|_{q} > \delta - 2\varepsilon_{i}\) .
\end{enumerate}

(Proceed by induction and use the formulas

\[
u = u _ {n, m} + p _ {n} u (1 - r _ {m}) + (1 - p _ {n}) u r _ {m} + p _ {n} u q _ {m}
\]

for n and m in N.)

\begin{enumerate}
\def\labelenumi{\alph{enumi})}
\setcounter{enumi}{2}
\item
  Deduce by contradiction that u is compact ("Pitt's theorem").
\item
  Show that the canonical injection of \(\ell_{\mathrm{K}}^{p}(\mathbf{N})\) into \(\ell_{\mathrm{K}}^{r}(\mathbf{N})\) is not compact.
\item
  Let p and r be real numbers such that \(1 \leqslant p \leqslant r \leqslant +\infty\) . Show that \(\ell_{\mathrm{K}}^{r}(\mathbf{N})\) is isomorphic to \(\ell_{\mathrm{K}}^{p}(\mathbf{N})\) if and only if p = r.
\end{enumerate}

\begin{enumerate}
\def\labelenumi{\arabic{enumi})}
\setcounter{enumi}{10}
\tightlist
\item
  Let E and F be Banach spaces over K. A linear map \(u \in \mathcal{L}(\mathrm{E};\mathrm{F})\) is said to be completely continuous if, for every weakly convergent sequence \((x_{n})_{n \in \mathbb{N}}\) in E, the sequence \((u(x_{n}))_{n \in \mathbb{N}}\) converges in F.
\end{enumerate}

\begin{enumerate}
\def\labelenumi{\alph{enumi})}
\item
  Every compact linear map is completely continuous.
\item
  If the dual space E' of E is countably generated, then every completely continuous linear map E → F is compact.
\item
  Let \((u_{n})_{n\in\mathbf{N}}\) be a sequence of elements of \(\ell_{\mathrm{K}}^{1}(\mathbf{N})\) . Show that \((u_{n})\) converges if and only if it converges weakly ("Schur's theorem"; cf. EVT, IV, p. 54, exercise 15).
\end{enumerate}

\(d)\) Deduce that the injection of \(\ell_{\mathrm{K}}^{1}(\mathbf{N})\) into \(\ell_{\mathrm{K}}^{2}(\mathbf{N})\) is completely continuous but is not compact.

\begin{enumerate}
\def\labelenumi{\alph{enumi})}
\setcounter{enumi}{4}
\item
  For every Banach space, every continuous linear map \(\mathrm{E} \to \ell_{\mathrm{K}}^{1}(\mathbf{N})\) is completely continuous, and every continuous linear map \(\ell_{\mathrm{K}}^{1}(\mathbf{N}) \to \mathrm{E}\) is completely continuous.
\item
  If E is reflexive or if its dual is countably generated, every continuous linear map \(\mathrm{E} \rightarrow \ell_{\mathrm{K}}^{1}(\mathbf{N})\) is compact.
\end{enumerate}

\begin{enumerate}
\def\labelenumi{\arabic{enumi})}
\setcounter{enumi}{11}
\tightlist
\item
  Let E be an infinite-dimensional reflexive Banach space. Let \(u \in \mathcal{L}^{c}(E)\) be a compact endomorphism. Suppose that u is not of finite rank. Let \(E_{\sigma}\) be the space E equipped with the weak topology.
\end{enumerate}

Show that the linear map \(^{t}u\) is a compact endomorphism of \((\mathrm{E}_{\sigma})'\) .

\begin{enumerate}
\def\labelenumi{\arabic{enumi})}
\setcounter{enumi}{12}
\tightlist
\item
  Let E be a countably generated Banach space over K. Recall (EVT, IV, p. 70, exercise 14) that a sequence \((x_{n})_{n\in \mathbf{N}}\) of elements of E is called a Schauder basis of E if, for every \(x\in \mathrm{E}\) , there exists a unique sequence of scalars \((\lambda_n(x))_{n\in \mathbf{N}}\) such that
\end{enumerate}

\[
x = \sum_ {n = 1} ^ {+ \infty} \lambda_ {n} (x) x _ {n}
\]

where the series converges in E.

\begin{enumerate}
\def\labelenumi{\alph{enumi})}
\item
  Every separable Hilbert space admits a Banach basis; every space \(\ell^{p}(\mathrm{I})\) , where I is countable, admits a Banach basis.
\item
  The Banach space E of bounded holomorphic functions on the open disk D of radius 1 in the complex plane, equipped with the norm \(\| f \| = \sup_{z \in \mathrm{D}} |f(z)|\) , has a Schauder basis.
\item
  Every Banach space E admitting a Banach basis is an approximation space. More precisely, there exists a real number \(c \geqslant 0\) such that \(1_{E}\) belongs to the closure in the space \(\mathcal{L}_{\mathrm{pc}}(\mathrm{E})\) of the set of finite-rank endomorphisms with norm \(\leqslant c\) in E.
\end{enumerate}

\begin{enumerate}
\def\labelenumi{\arabic{enumi})}
\setcounter{enumi}{13}
\tightlist
\item
  \begin{enumerate}
  \def\labelenumii{\alph{enumii})}
  \tightlist
  \item
    Let E and F be Hilbert spaces. A continuous linear map u from E into F is compact if and only if the image of u contains no closed subspace of infinite dimension.
  \end{enumerate}
\end{enumerate}

\begin{enumerate}
\def\labelenumi{\alph{enumi})}
\setcounter{enumi}{1}
\tightlist
\item
  The characterization in a) of compact endomorphisms from E into F does not hold for every pair of Banach spaces (E, F). (One may use Exercise 5, p. 121.)
\end{enumerate}

\begin{enumerate}
\def\labelenumi{\arabic{enumi})}
\setcounter{enumi}{14}
\tightlist
\item
  Let \(E_{0}\) be the complex vector space of functions \(f: N^{*} \times N^{*} \to C\) . For \(k \in N^{*}\) , let \(\alpha_{k} \in E_{0}\) be the function such that
\end{enumerate}

\[
\alpha_ {k} (n, m) = \left\{ \begin{array}{l l} m ^ {k} & \text { if } 1 \leqslant n \leqslant k - 1 \\ n ^ {k} & \text { if } n \geqslant k. \end{array} \right.
\]

Let E be the vector subspace of \(E_{0}\) formed by the functions f such that \(p_{k}(f)\) is finite for all \(k \in N^{*}\) , where \(p_{k}\) is defined by

\[
p _ {k} (f) = \sum_ {n, m} \alpha_ {k} (n, m) | f (n, m) |.
\]

\begin{enumerate}
\def\labelenumi{\alph{enumi})}
\tightlist
\item
  The space E equipped with the topology defined by the norms \((p_{k})_{k\in\mathbf{N}^{*}}\) is a Fréchet space and a Montel space. Its dual E' is identified with the subspace of \(E_{0}\) consisting of functions g such that, for at least one \(k\geqslant1\) , one has
\end{enumerate}

\[
\sup _ {m, n} \alpha_ {k} (m, n) ^ {- 1} | g (m, n) | <   + \infty
\]

(cf. EVT, IV, p. 63, exercise 8).

\begin{enumerate}
\def\labelenumi{\alph{enumi})}
\setcounter{enumi}{1}
\tightlist
\item
  Consider the linear map \(u\) from E into \(\ell^1 (\mathbf{N}^*)\) defined by
\end{enumerate}

\[
u (f) = \left(\sum_ {n} f (n, m)\right) _ {m \in \mathbf {N} ^ {*}}.
\]

The linear map u is continuous.

\begin{enumerate}
\def\labelenumi{\alph{enumi})}
\setcounter{enumi}{2}
\item
  Identify \(\ell^{\infty}(\mathbf{N}^{*})\) with the dual of \(\ell^{1}(\mathbf{N}^{*})\) . The transpose of \(u\) is the map \(v\colon \ell^{\infty}(\mathbf{N}^{*})\to \mathrm{E}^{\prime}\subset \mathrm{E}_{0}\) determined by \((x_{n})\mapsto f\) where \(f(n,m) = x_{m}\) for all \((n,m)\in \mathbf{N}^{*}\times \mathbf{N}^{*}\) .
\item
  The map u induces an isomorphism \(\mathrm{E}/\mathrm{Ker}(u)\to\ell^{1}(\mathbf{N}^{*})\) (cf. loc. cit.).
\item
  The map u is not compact.
\item
  The transpose v of u is compact for the topology of convergence on bounded subsets of \(\ell^{\infty}(\mathbf{N}^{*})\) .
\item
  The map v is an example of a compact linear map between separated locally convex spaces whose transpose is not compact.
\end{enumerate}

\begin{enumerate}
\def\labelenumi{\arabic{enumi})}
\setcounter{enumi}{15}
\item
  Let E be a locally convex space and u a compact endomorphism of E. Denote by \(E_{b}^{\prime}\) the dual space of E equipped with the topology of bounded convergence. The linear map \(^{t}u \circ ^{t}u\) is a compact endomorphism of \(E_{b}^{\prime}\) .
\item
  Let E be a Banach space of countable type. Then \(\mathcal{L}^{\mathrm{c}}(\mathrm{E})\) , equipped with the topology induced by the norm of \(\mathcal{L}(\mathrm{E})\) is of countable type if, and only if, the dual space \(\mathrm{E}'\) is of countable type.
\item
  Let E, F, G be Banach spaces and let u be an element of \(\mathcal{L}^{c}(\mathrm{E};\mathrm{F})\) .
\end{enumerate}

\begin{enumerate}
\def\labelenumi{\alph{enumi})}
\tightlist
\item
  Let v be an element of \(\mathcal{L}(\mathrm{E};\mathrm{G})\) . Suppose that, for every \(\varepsilon > 0\) , there exists C \textgreater{} 0 such that
\end{enumerate}

\[
\| v (x) \| _ {\mathrm{G}} \leqslant \varepsilon \| x \| _ {\mathrm{E}} + \mathrm{C} \| u (x) \| _ {\mathrm{F}}
\]

for all \(x\in \mathrm{E}\) . Then \(v\in \mathcal{L}^{\mathrm{c}}(\mathrm{E};\mathrm{G})\) .

\begin{enumerate}
\def\labelenumi{\alph{enumi})}
\setcounter{enumi}{1}
\tightlist
\item
  Let w be an injective element of \(\mathcal{L}(\mathrm{F};\mathrm{G})\) . Then, for every \(\varepsilon > 0\) , there exists C \textgreater{} 0 such that
\end{enumerate}

\[
\| u (x) \| _ {\mathrm{F}} \leqslant \varepsilon \| x \| _ {\mathrm{E}} + \mathrm{C} \| (w \circ u) (x) \| _ {\mathrm{G}}
\]

for all \(x \in E\) .

\begin{enumerate}
\def\labelenumi{\arabic{enumi})}
\setcounter{enumi}{18}
\tightlist
\item
  Let X be a locally compact, metrizable, and countable at infinity topological space, and \(\mu\) a positive measure on X. We denote by A the set of atoms of \(\mu\) . Let \(f \in \mathcal{L}^{\infty}(X, \mu)\) and \(p \in [1, +\infty]\) .
\end{enumerate}

Show that multiplication by \(f\) defines a compact endomorphism of \(\mathrm{L}^{p}(\mathrm{X},\mu)\) if and only if the following two conditions are satisfied:

\begin{enumerate}
\def\labelenumi{(\roman{enumi})}
\item
  The set \(\{x\in X-A\mid f(x)\neq0\}\) is \(\mu\) -negligible;
\item
  For every \(\varepsilon > 0\) there exists a finite subset F of A such that \(|f(a)| \leqslant \varepsilon\) for all \(a \in \mathbf{A} - \mathbf{F}\) .
\end{enumerate}

(Consider multiplication by \(f\) on the space \(\mathrm{L}^{p}(\mathrm{X}_{\varepsilon},\mu)\) , where \(\mathrm{X}_{\varepsilon}\) is the set of \(x\in\mathrm{X}\) such that \(|f(x)|\geqslant\varepsilon\) .)

\begin{enumerate}
\def\labelenumi{\arabic{enumi})}
\setcounter{enumi}{19}
\tightlist
\item
  Let X, Y be compact metric spaces, T a compact operator from \(\mathcal{C}(\mathrm{Y})\) into \(\mathcal{C}(\mathrm{X})\) .
\end{enumerate}

Prove that there exists a positive measure \(\mu\) on \(X\) and a measurable map \(N\) from \(X \times Y\) into \(\mathbf{C}\) such that

\begin{enumerate}
\def\labelenumi{\alph{enumi})}
\item
  For every \(y \in Y\) , the map \(\mathrm{N}_{y} \colon x \mapsto \mathrm{N}(x, y)\) from \(X\) into \(\mathbf{C}\) is \(\mu\) -integrable;
\item
  The map \(y \mapsto N_{y}\) from Y into \(\mathcal{L}^{1}(X, \mu)\) is continuous;
\item
  For every \(f \in \mathcal{C}(\mathrm{X})\) , and every \(y \in \mathrm{Y}\) , we have
\end{enumerate}

\[
\mathrm{T} f (y) = \int_ {\mathrm{X}} \mathrm{N} (x, y) f (x) d \mu (x).
\]

\begin{enumerate}
\def\labelenumi{\arabic{enumi})}
\setcounter{enumi}{20}
\tightlist
\item
  We denote by D the open disk with center 0 and radius 1 in C, and by H \(^{2}\) (D) the Hardy space on D (EVT, V, p. 7, example 5). We write T = R/2πZ, which we equip with the normalized Haar measure. Let j be the isometry from H \(^{2}\) (D) into L \(^{2}\) (T) defined by
\end{enumerate}

\[
j (f) (\theta) = \sum_ {n \geqslant 0} \frac {f ^ {(n)} (0)}{n !} e ^ {i n \theta}.
\]

The image of H \(^{2}\) (D) by \(j\) will be denoted H \(^{2}\) (T).

Let \(\widetilde{\mathrm{P}}\) be the orthogonal projector of \(\mathrm{L}^2 (\mathbf{T})\) with image \(\mathrm{H}^2 (\mathbf{T})\), and let P be the linear map \(f\mapsto \widetilde{\mathrm{P}} (f)\) from \(\mathrm{L}^2 (\mathbf{T})\) into \(\mathrm{H}^2 (\mathbf{T})\) . If \(f\) is an element of \(\mathrm{L}^{\infty}(\mathbf{T})\), let \(\mathrm{M}_f\) be the endomorphism \(g\mapsto fg\) of \(\mathrm{L}^2 (\mathbf{T})\) .

Show that if \(f\) is a continuous function, then the commutator \(\widetilde{\mathrm{P}}\mathrm{M}_f - \mathrm{M}_f\widetilde{\mathrm{P}}\) is a compact linear map. (Consider the particular case where \(f\) is a trigonometric polynomial.)

\begin{enumerate}
\def\labelenumi{\arabic{enumi})}
\setcounter{enumi}{21}
\tightlist
\item
  We keep the notations from the previous exercise. Let R be the continuous linear map from \(\mathrm{H}^{2}(\mathbf{T})\) into \(\mathrm{L}^{2}(\mathbf{T})\) defined by \(\mathrm{R}f(\theta)=f(-\theta)\) . For \(g\in\mathrm{L}^{\infty}(\mathbf{T})\) we set \(\Gamma_{g}=\mathrm{PM}_{g}\mathrm{R}\) .
\end{enumerate}

\begin{enumerate}
\def\labelenumi{\alph{enumi})}
\item
  Prove that the linear map \(\Gamma\colon g\mapsto\Gamma_{g}\) from \(\mathrm{L}^{\infty}(\mathbf{T})\) into \(\mathcal{L}(\mathrm{H}^{2}(\mathbf{T}))\) is continuous, with norm \(\leqslant 1\), and its kernel F is equal to \(\mathrm{H}^{2}(\mathbf{T})^{\circ}\cap\mathrm{L}^{\infty}(\mathbf{T})\) .
\item
  If \(g\) is a continuous function, then \(\Gamma_g\) is a compact endomorphism of \(\mathrm{H}^2 (\mathbf{T})\) .
\item
  Let \(g \in \mathrm{L}^{\infty}(\mathbf{T})\) . We denote by \(\alpha_{g}\) the linear form on \(\mathrm{H}^{2}(\mathbf{T})\) defined by
\end{enumerate}

\[
\alpha_ {g} (f) = \int_ {\mathbf {T}} \Gamma_ {g} (f) (\theta) d \theta .
\]

Let \(\mathcal{A}\) be the image under \(j\) of the subspace of \(\mathrm{H}^2 (\mathrm{D})\) consisting of holomorphic functions that extend analytically to a neighborhood of \(\overline{\mathrm{D}}\) . Prove that if \(f_{1}\) and \(f_{2}\) belong to \(\mathcal{A}\) , then

\[
\alpha_ {g} (f _ {1} f _ {2}) = \int_ {\mathbf {T}} \Gamma_ {g} (f _ {1}) (\theta) f _ {2} (- \theta) d \theta .
\]

\(d)\) Let \(\varphi\) be a holomorphic function in a neighborhood of \(\overline{\mathrm{D}}\) not vanishing on \(\partial\mathrm{D}\) and let \(f = j(\varphi|\mathrm{D})\) . Then there exist \(f_{1}\) and \(f_{2}\) in \(\mathscr{A}\) such that \(f = f_{1}f_{2}\) and

\[
\| f \| _ {\mathrm{L} ^ {1} (\mathbf {T})} = \| f _ {1} \| _ {\mathrm{L} ^ {2} (\mathbf {T})} \| f _ {2} \| _ {\mathrm{L} ^ {2} (\mathbf {T})}.
\]

(Let \(z_{1},\ldots,z_{n}\) be the zeros of \(\varphi\) in D; then there exists a function \(\psi\) holomorphic in a neighborhood of D such that

\[
\varphi (z) = \prod_ {k = 1} ^ {n} \frac {z - z _ {k}}{1 - \overline {{z}} _ {k} z} \psi (z) ^ {2}
\]

for all z.)

\begin{enumerate}
\def\labelenumi{\alph{enumi})}
\setcounter{enumi}{4}
\tightlist
\item
  Prove that, for all \(f \in \mathrm{H}^{2}(\mathbf{T})\) , we have the inequality
\end{enumerate}

\[
\left| \alpha_ {g} (f) \right| \leqslant \left\| \Gamma_ {g} \right\| _ {\mathscr {L} \left(\mathrm{H} ^ {2} (\mathbf {T})\right)} \| f \| _ {\mathrm{L} ^ {1} (\mathbf {T})}.
\]

(If \(\varphi \in \mathrm{H}^2 (\mathrm{D})\) , there exists a sequence \((r_n)_{n\in \mathbf{N}}\) of ]0,1[ such that, for every \(n\in \mathbf{N}\) , the function \(\varphi_{n}\) defined by \(\varphi_{n}(z) = \varphi (r_{n}z)\) has no zeros on \(\partial \mathrm{D}\) , and \(\lim \varphi_n = \varphi\) in \(\mathrm{H}^2 (\mathrm{D}).\)

Deduce that there exists \(h \in \mathrm{L}^{\infty}(\mathbf{T})\) such that

\[
\alpha_ {g} (f) = \int_ {\mathbf {T}} f (\theta) h (\theta) d \theta
\]

for all \(f \in \mathrm{H}^2(\mathbf{T})\) , and \(\|h\|_{\mathrm{L}^\infty(\mathbf{T})} \leqslant \| \Gamma_g \|_{\mathcal{L}(\mathrm{H}^2(\mathbf{T}))}\) .

\(f)\) Prove that there exists \(k \in \mathrm{L}^{\infty}(\mathbf{T})\) such that \(\Gamma_g = \Gamma_k\) and \(\| k \|_{\mathrm{L}^{\infty}(\mathbf{T})} = \| \Gamma_g \|_{\mathcal{L}(\mathrm{H}^2(\mathbf{T}))}\) .

The map \(\widehat{\Gamma}\colon\mathrm{L}^{\infty}(\mathbf{T})/\mathrm{F}\to\mathcal{L}(\mathrm{H}^{2}(\mathbf{T}))\) is therefore an isometry ("Nehari's theorem").

\(g)\) Let \(g\) be a continuous function on \(\mathbf{T}\) . Then we have \(d(g, \mathrm{F}) = d(g, \mathrm{F} \cap \mathcal{C}(\mathbf{T}))\) . (For \(\mathrm{R} > 1\) , one defines the linear map \(\mathrm{P}_{\mathrm{R}}\) from \(\mathrm{L}^{\infty}(\mathbf{T})\) into \(\mathcal{C}(\mathbf{T})\) by

\[
\mathrm{P} _ {\mathrm{R}} f (\theta) = \int_ {\mathbf {T}} \frac {\mathrm{R} ^ {2} - 1}{| \mathrm{R} e ^ {i (\theta - \varphi)} - 1 | ^ {2}} f (\varphi) d \varphi .
\]

Then \(\| \mathrm{P}_{\mathrm{R}}f\|_{\mathcal{C}(\mathbf{T})}\leqslant \| f\|_{\mathrm{L}^{\infty}(\mathbf{T})}\) and, if \(f\in \mathcal{C}(\mathbf{T})\) , we have \(\mathrm{P_R}f\to f\) in \(\mathcal{C}(\mathbf{T})\) as \(\mathrm{R}\rightarrow 1.\)

Deduce that \(\Gamma(\mathcal{C}(\mathbf{T}))\) is a closed subspace of \(\mathcal{L}(\mathrm{H}^{2}(\mathbf{T}))\) ("Sarason's theorem").

\begin{enumerate}
\def\labelenumi{\alph{enumi})}
\setcounter{enumi}{7}
\tightlist
\item
  Let \(g \in \mathrm{L}^{\infty}(\mathbf{T})\) such that \(\Gamma_{g}\) is a compact endomorphism. For \(m \in Z\) we set \(e_{m}(\theta) = e^{im\theta}\) . Prove that \(\|\Gamma_{ge_{-m}}\|_{\mathcal{L}(\mathrm{H}^{2}(\mathbf{T}))}\) tends to zero as m tends to \(+\infty\) , and then that there exists a sequence \((h_{m})_{m \in \mathbf{N}}\) of F such that the sequence \((e_{m}h_{m})_{m \in \mathbf{N}}\) converges to g in \(\mathrm{L}^{\infty}(\mathbf{T})\) . Deduce that there exists a continuous function in the class of g modulo F ("Hartman's theorem").
\end{enumerate}

\begin{enumerate}
\def\labelenumi{\arabic{enumi})}
\setcounter{enumi}{22}
\item
  Let B be an infinite-dimensional Banach space. We denote by E the strong dual of B and by F the dual of B equipped with the topology \(\sigma(\mathrm{B}^{\prime},\mathrm{B})\) . Prove that the identity map of the dual of B is compact from E into F, but that its transpose is not compact from \(\mathrm{F}_{b}^{\prime}\) into \(\mathrm{E}_{b}^{\prime}\) .
\item
  Let E be a real or complex Banach space and \(u \in \mathcal{L}(\mathrm{E})\) . Assume that the sequence \((u^{n})_{n \in \mathbf{N}}\) is bounded and that u is compact from E into E equipped with the weak topology \(\sigma(\mathrm{E}, \mathrm{E}')\) . For every integer \(N \geqslant 1\) we set
\end{enumerate}

\[
v _ {\mathrm{N}} = \frac {1}{\mathrm{N}} (1 _ {\mathrm{E}} + u + \dots + u ^ {\mathrm{N-1}}) \in \mathcal {L} (\mathrm{E}).
\]

\begin{enumerate}
\def\labelenumi{\alph{enumi})}
\item
  Let F be the closure of the image of \(1_{\mathrm{E}} - u\) . For every \(x \in \mathrm{F}\) , we have \(v_{\mathrm{N}}(x) \to 0\) in E when \(\mathrm{N} \to +\infty\) .
\item
  Let \(x \in E\) . The sequence \((v_{\mathrm{N}}(x))\) converges to an element \(y \in E\) such that \(u(y) = y\) . (There exists a subsequence of \((v_{\mathrm{N}}(x))\) that converges weakly to \(y \in E\) ; show that \(x - y \in F\) using the Hahn--Banach theorem.)
\item
  The map \(w \colon \mathrm{E} \to \mathrm{E}\) such that \(w(x)\) is the element \(y\) given by the preceding question is a continuous projector whose image is the eigenspace of \(u\) corresponding to the eigenvalue 1.
\end{enumerate}

\(d)\) Let \(z \in \mathbf{C}\) such that \(|z| = 1\) . The sequence in \(\mathcal{L}(\mathrm{E})\) defined by

\[
v _ {z, \mathrm{N}} = \frac {1}{\mathrm{N}} (1 _ {\mathrm{E}} + z ^ {- 1} u \dots + z ^ {- (\mathrm{N-1})} u ^ {\mathrm{N-1}})
\]

converges in \(\mathcal{L}(\mathrm{E})\) equipped with the topology of pointwise convergence. Let \(w_{z}\) be its limit. Then \(w_{z}\) is a projector, one has \(w_{z}w = ww_{z} = zw_{z}\) and \(w_{z_1}w_{z_2} = 0\) if \(z_{1} \neq z_{2}\) .

\begin{enumerate}
\def\labelenumi{\arabic{enumi})}
\setcounter{enumi}{24}
\tightlist
\item
  \begin{enumerate}
  \def\labelenumii{\alph{enumii})}
  \tightlist
  \item
    Let \(N \geqslant 1\) be an integer. We denote by \(G_{N}\) (resp. \(H_{N}\) ) the set \(\{-1,1\}^{N}\) (resp. the set \(\{-1,2\}^{N}\) ). We denote by \(\mu_{N}\) (resp. \(\nu_{N}\) ) the positive measure of total mass 1 on \(G_{N}\) (resp. on \(H_{N}\) ) such that, for every \((t_{i})_{1 \leqslant i \leqslant N}\) in \(G_{N}\) (resp. in \(H_{N}\) ), one has \(\mu_{\mathrm{N}}((t_{i})_{1 \leqslant i \leqslant N}) = \frac{1}{2^{\mathrm{N}}}\) , resp.
  \end{enumerate}
\end{enumerate}

\[
\nu_ {\mathrm{N}} ((t _ {i}) _ {1 \leqslant i \leqslant \mathrm{N}}) = \left(\frac {1}{3}\right) ^ {j} \left(\frac {2}{3}\right) ^ {\mathrm{N} - j},
\]

where j is the number of indices i such that \(t_{i}=2\) .

There exists a real number \(c > 0\) such that, for every family \((a_{i})_{1\leqslant i\leqslant N}\) of complex numbers, we have

\[
\mu_ {\mathrm{N}} \Big (\Big \{(t _ {i}) _ {1 \leqslant i \leqslant \mathrm{N}} \in \mathrm{G} _ {\mathrm{N}} | \Big | \sum_ {i = 1} ^ {\mathrm{N}} a _ {i} t _ {i} \Big | > c (\log \mathrm{N}) ^ {1 / 2} \Big (\sum_ {i = 1} ^ {\mathrm{N}} | a _ {i} | ^ {2} \Big) ^ {1 / 2} \Big \} \Big) <   \frac {c}{\mathrm{N} ^ {3}},
\]

and

\[
\nu_ {\mathrm{N}} \Big (\Big \{(t _ {i}) _ {1 \leqslant i \leqslant \mathrm{N}} \in \mathrm{H} _ {\mathrm{N}} | \left| \sum_ {i = 1} ^ {\mathrm{N}} a _ {i} t _ {i} \right| > c (\log \mathrm{N}) ^ {1 / 2} \Big (\sum_ {i = 1} ^ {\mathrm{N}} | a _ {i} | ^ {2} \Big) ^ {1 / 2} \Big \} \Big) <   \frac {c}{\mathrm{N} ^ {3}}.
\]

(We may assume that \(a_{i} \in \mathbf{R}\) for all \(i\) and that \(\sum |a_{i}|^{2} = 1\) . For the first inequality, prove and use the upper bound.

\[
\int_ {G _ {N}} \exp \left(\lambda \left| \sum_ {i = 1} ^ {N} a _ {i} t _ {i} \right|\right) d \mu_ {N} \leqslant 2 \exp (\lambda^ {2})
\]

for every real \(\lambda > 0\) ; reason in a similar manner for the second upper bound using the inequality \(\frac{1}{3}(e^{2x} + 2e^{-x}) \leqslant e^{2x^2}\) for all \(x \in \mathbf{R}\) .)

\begin{enumerate}
\def\labelenumi{\alph{enumi})}
\setcounter{enumi}{1}
\tightlist
\item
  For \(k \in N\) , let \(\mathrm{A}_{k} = \mathbf{Z}/(3 \cdot 2^{k})\mathbf{Z}\) . There exists a real number \(c' > 0\) , not depending on k, and a function \(f_{k} \colon \widehat{\mathrm{A}}_{k} \to \{-1, 2\}\) such that
\end{enumerate}

\[
\sum_ {\chi \in \widehat {\mathrm{A}} _ {k}} f _ {k} (\chi) = 0
\]

and, for all \(x \in A_{k}\) , one has

\[
\Big | \sum_ {\chi \in \widehat {\mathrm{A}} _ {k}} \chi (x) f _ {k} (\chi) \Big | \leqslant c ^ {\prime} (k + 1) ^ {1 / 2} 2 ^ {k / 2}.
\]

(Using \(a\) ), determine the existence of a function \(f_{k}\) satisfying the second condition, then modify \(f_{k}\) to obtain the first.)

We denote by \(B_{k}\) (resp. \(C_{k}\) ) the set of characters \(\chi\) of \(A_{k}\) such that \(f_{k}(\chi)=2\) (resp. \(f_{k}(\chi)=-1\) ). We have \(\mathrm{Card(B_{k})}=2^{k}\) and \(\mathrm{Card(C_{k})}=2^{k+1}\) .

For every \(k \geqslant 1\) , we fix a bijection \(\iota_k \colon B_k \to C_{k-1}\) .

We denote by A the disjoint union of the \(A_{k}\) for \(k \in N\) ; we equip A with the discrete topology, and we denote by \(\mathrm{E} = \mathcal{C}_{b}(\mathrm{A})\) .

We fix a family \(\varepsilon = (\varepsilon_k)_{k\geqslant \mathbf{N}}\) of functions \(\varepsilon_{k}\colon \mathrm{B}_{k}\to \{-1,1\}\) . We then define functions \(e_{k,\chi}\in \mathrm{E}\) , for \(k\in \mathbf{N}\) and \(\chi \in \mathrm{B}_k\) by

\[
e _ {k, \chi} (x) = \left\{ \begin{array}{l l} \iota_ {k} (\chi) (x) & \text { if } k \geqslant 1 \text {   and   } x \in \mathrm{A} _ {k - 1}, \\ \varepsilon_ {k} (\chi) \chi (x) & \text { if } x \in \mathrm{A} _ {k} \\ 0 & \text { otherwise. } \end{array} \right.
\]

We denote by \(E_{\varepsilon}\) the closed subspace of E generated by the functions \((e_{k,\chi})_{k\in\mathbf{N},\chi\in\mathbf{B}_{k}}\) . This is a complex Banach space of countable type. The rest of the exercise aims to show that, for a suitable choice of \(\varepsilon\) , the space \(\mathrm{E}_{\varepsilon}\) does not have the approximation property; it admits no Banach basis.

\begin{enumerate}
\def\labelenumi{\alph{enumi})}
\setcounter{enumi}{2}
\tightlist
\item
  For \(k \in N\) and \(\chi \in B_{k}\) there exists a continuous linear form \(\lambda_{k,\chi} \in E_{\varepsilon}'\) such that, for every \(f \in E_{\varepsilon}\) ,
\end{enumerate}

\[
\lambda_ {k, \chi} (f) = \frac {1}{\operatorname{Card} (\mathrm{A} _ {k})} \sum_ {x \in \mathrm{A} _ {k}} \varepsilon_ {k} (\chi) \chi (x ^ {- 1}) f (x).
\]

For \(k \geqslant 1\) and \(f \in \mathrm{E}_{\varepsilon}\) , we have

\[
\lambda_ {k, \chi} (f) = \frac {2}{\operatorname{Card} (\mathrm{A} _ {k})} \sum_ {x \in \mathrm{A} _ {k}} \iota_ {k} (\chi) (x ^ {- 1}) f (x).
\]

\begin{enumerate}
\def\labelenumi{\alph{enumi})}
\setcounter{enumi}{3}
\tightlist
\item
  For every \(k \in N\) , let \(\beta_{k}\) be the continuous linear form on \(\mathcal{L}(E_{\varepsilon})\) defined by
\end{enumerate}

\[
\beta_ {k} (u) = \frac {1}{2 ^ {k}} \sum_ {\chi \in \mathrm{B} _ {k}} \lambda_ {k, \chi} (u (e _ {k, \chi})).
\]

One has \(\beta_{k}(1_{\mathrm{E}_{\varepsilon}}) = 1\)

\begin{enumerate}
\def\labelenumi{\alph{enumi})}
\setcounter{enumi}{4}
\tightlist
\item
  Let \(u \in \mathcal{L}(\mathrm{E}_{\varepsilon})\) and \(k \in \mathbf{N}\) . We have
\end{enumerate}

\[
\beta_ {k + 1} (u) - \beta_ {k} (u) = \frac {1}{\operatorname{Card} \left(\mathrm{A} _ {k}\right)} \sum_ {x \in \mathrm{A} _ {k}} u \left(\varphi_ {k, x}\right) (x)
\]

where, for \(x \in \mathrm{A}_k\) , we set

\[
\varphi_ {k, x} = \frac {1}{2 ^ {k + 1}} \sum_ {\chi \in B _ {k + 1}} \iota_ {k + 1} (\chi) (x ^ {- 1}) e _ {k + 1, \chi} - \frac {1}{2 ^ {k}} \sum_ {\chi \in B _ {k}} \varepsilon_ {k} (\chi) \chi (x ^ {- 1}) e _ {k, \chi}.
\]

One has \(|\varphi_{k,x}(y)| \leqslant c'(k + 1)^{1/2}2^{-k/2}\) for \((x,y) \in \mathrm{A}_k^2\) .

\begin{enumerate}
\def\labelenumi{\alph{enumi})}
\setcounter{enumi}{5}
\item
  There exists a real number \(c'' \geqslant c'\) such that, for every \(k \geqslant 1\) , there exists a function \(\varepsilon_k \colon B_k \to \{-1,1\}\) satisfying \(|\varphi_{k,x}(y)| \leqslant c''3(k+1)^{1/2}2^{-k/2}\) for \((x,y) \in A_k \times A_{k-1}\) . (Apply the first inequality of a).) We now assume that \(\varepsilon = (\varepsilon_k)\) satisfies these properties.
\item
  For every \(k \in \mathbf{N}\) and \(x \in \mathrm{A}_k\) , we have \(\| \varphi_{k,x} \| \leqslant c''(k + 1)^{1/2} 2^{-k/2}\) .
\item
  Let K be the subset of \(E_{\varepsilon}\) whose elements are \(e_{0,\chi_{0}}\) , where \(\chi_{0}\) is the unique element of \(B_{0}\) , and the functions \((k+1)^{2}\varphi_{k,x}\) for \(k\in N\) and \(x\in A_{k}\) . The set K is relatively compact in \(E_{\varepsilon}\) , and we have
\end{enumerate}

\[
| \beta_ {k + 1} (u) - \beta_ {k} (u) | \leqslant \frac {1}{(k + 1) ^ {2}} \sup _ {x \in K} \| u (x) \|
\]

for all \(k \in \mathbf{N}\) and every \(u \in \mathcal{L}(\mathrm{E}_{\varepsilon})\) .

\begin{enumerate}
\def\labelenumi{\roman{enumi})}
\tightlist
\item
  The sequence \((\beta_{k})_{k\in \mathbf{N}}\) converges in the dual of the space \(\mathcal{L}(\mathrm{E}_{\varepsilon})\) equipped with the topology of pointwise convergence. Denote by \(\beta\) its limit. For every \(u\in \mathcal{L}(\mathrm{E}_{\varepsilon})\) , we have \(|\beta (u)|\leqslant 3\sup_{x\in \mathrm{K}}\| u(x)\|\) .
\end{enumerate}

\begin{enumerate}
\def\labelenumi{\alph{enumi})}
\setcounter{enumi}{9}
\tightlist
\item
  One has \(\beta(1_{\mathrm{E}_{\varepsilon}})=1\) and \(\beta(u)=0\) if \(u\) has finite rank.
\end{enumerate}

\(k)\) The space \(\mathrm{E}_{\varepsilon}\) does not have the approximation property.

\begin{enumerate}
\def\labelenumi{\alph{enumi})}
\setcounter{enumi}{11}
\tightlist
\item
  Let \(p \in]2, +\infty[\) . We equip A with the counting measure and denote by \(E_{\varepsilon}^{p}\) the closed subspace of \(L^{p}(A)\) generated by the \((e_{k,\chi})_{k \in \mathbf{N}, \chi \in \mathrm{B}_{k}}\) . The space \(E_{\varepsilon}^{p}\) does not have the approximation property. (Verify that the cardinality of the support of \(\varphi_{k,x}\) is \(\leqslant 21 \cdot 2^{k}\) and deduce that \(\|\varphi_{k,x}\|_{p} = \mathrm{O}((k+1)^{1/2}2^{-k(p-2)/2p})\) for \(k \in N\) and \(x \in A_{k}\) .)
\end{enumerate}

\begin{enumerate}
\def\labelenumi{\arabic{enumi})}
\setcounter{enumi}{25}
\tightlist
\item
  Let E and F be Banach spaces, and let u be a continuous linear map from E into F.
\end{enumerate}

\begin{enumerate}
\def\labelenumi{\alph{enumi})}
\tightlist
\item
  Prove that the following conditions are equivalent:
\end{enumerate}

\begin{enumerate}
\def\labelenumi{(\roman{enumi})}
\item
  The map u is compact;
\item
  There exists a sequence \((x_n')_n\) tending to 0 in the strong dual \(\mathrm{E}'\) of E satisfying \(\| u(x)\| \leqslant \sup |\langle x,x_n' \rangle|\) for all \(x \in \mathrm{E}\) ;
\item
  There exists a closed subspace L of the Banach space \(c_{0}\) and compact linear maps \(f: E \to L\) and \(g: L \to F\) such that \(u = g \circ f\) (If u is compact, apply cor. 1 of EVT, IV, p. 24 to the image under \(^{t}u\) of the polar of the unit ball of F. Under hypothesis (ii), let \((\lambda_{n})\) be an element of \(c_{0}\) such that the sequence \((\lambda_{n}^{-1}x_{n}^{\prime})\) tends to 0; define f by \(f(x) = (\langle x, \lambda_{n}^{-1}x_{n}^{\prime} \rangle)_{n}\) and take for L the closure of the image of f.)
\end{enumerate}

\begin{enumerate}
\def\labelenumi{\alph{enumi})}
\setcounter{enumi}{1}
\tightlist
\item
  If \(u\) is compact, we have \(\|u\| = \inf(\sup_n \|x_n'\|)\) where the infimum is taken over the set of sequences \((x_n')\) satisfying (ii), and \(\|u\| = \inf(\|f\|\|g\|)\) , the infimum being taken over the set of triples (L, \(f,g\) ) satisfying the conditions of (iii).
\end{enumerate}

\begin{enumerate}
\def\labelenumi{\arabic{enumi})}
\setcounter{enumi}{26}
\tightlist
\item
  Let E and F be locally convex spaces. A mapping \(u \in \mathcal{L}(\mathrm{E};\mathrm{F})\) is said to be weakly compact if it is a compact mapping from E into F equipped with the topology \(\sigma(\mathrm{F},\mathrm{F}')\) .
\end{enumerate}

\begin{enumerate}
\def\labelenumi{\alph{enumi})}
\item
  The weakly compact mappings form a subspace \(\mathcal{L}^{\mathrm{fc}}(\mathrm{E};\mathrm{F})\) of \(\mathcal{L}(\mathrm{E};\mathrm{F})\) . Let \(\mathrm{E}_1\) , \(\mathrm{F}_1\) be locally convex spaces, \(f\in \mathcal{L}(\mathrm{E}_1;\mathrm{E})\) , \(g\in \mathcal{L}(\mathrm{F};\mathrm{F}_1)\) and \(u\in \mathcal{L}(\mathrm{E};\mathrm{F})\) ; if \(u\) is weakly compact, the same is true of \(g\circ u\circ f\) .
\item
  One has \(\mathcal{L}^{c}(\mathrm{E};\mathrm{F})\subset\mathcal{L}^{\mathrm{fc}}(\mathrm{E};\mathrm{F})\) . If E is an infinite-dimensional reflexive Banach space, the identity map of E is weakly compact but is not compact.
\end{enumerate}

\begin{enumerate}
\def\labelenumi{\arabic{enumi})}
\setcounter{enumi}{27}
\tightlist
\item
  Let E and F be locally convex spaces and \(u \in \mathcal{L}(\mathrm{E};\mathrm{F})\) .
\end{enumerate}

\begin{enumerate}
\def\labelenumi{\alph{enumi})}
\tightlist
\item
  Show that the following properties are equivalent:
\end{enumerate}

\begin{enumerate}
\def\labelenumi{(\roman{enumi})}
\item
  The image under \(u\) of every bounded subset of E is relatively compact for the weak topology;
\item
  The map \(^{tt}u\) maps F'' into E.
\end{enumerate}

\begin{enumerate}
\def\labelenumi{\alph{enumi})}
\setcounter{enumi}{1}
\tightlist
\item
  If F is quasi-complete, prove that these properties are equivalent to:
\end{enumerate}

\begin{enumerate}
\def\labelenumi{(\roman{enumi})}
\setcounter{enumi}{2}
\tightlist
\item
  The map \(^{t}u\) maps every equicontinuous subset of F' to a relatively compact subset of E' for \(\sigma(\mathrm{E}',\mathrm{E}'')\) .
\end{enumerate}

\begin{enumerate}
\def\labelenumi{\alph{enumi})}
\setcounter{enumi}{2}
\tightlist
\item
  If E and F are Banach spaces, (i) to (iii) are also equivalent to
\end{enumerate}

(i') The map u is weakly compact;

(iii') The map \(^{t}u\colon F_{b}^{\prime}\to E_{b}^{\prime}\) is weakly compact.

\begin{enumerate}
\def\labelenumi{\arabic{enumi})}
\setcounter{enumi}{28}
\tightlist
\item
  \begin{enumerate}
  \def\labelenumii{\alph{enumii})}
  \tightlist
  \item
    Prove that every weakly compact subset of the Banach space \(\ell^{1}(\mathbf{N})\) is compact (cf. EVT, IV, p. 69, exerc. 11). In particular, every weakly compact linear map from a locally convex space E into \(\ell^{1}(\mathbf{N})\) is compact.
  \end{enumerate}
\end{enumerate}

\begin{enumerate}
\def\labelenumi{\alph{enumi})}
\setcounter{enumi}{1}
\tightlist
\item
  Let E be a Banach space. Deduce from a) that every weakly compact linear map from \(c_{0}\) into E is compact.
\end{enumerate}

\begin{enumerate}
\def\labelenumi{\arabic{enumi})}
\setcounter{enumi}{29}
\tightlist
\item
  \begin{enumerate}
  \def\labelenumii{\alph{enumii})}
  \tightlist
  \item
    Let E and F be Banach spaces and \(u \in \mathcal{L}(\mathrm{E};\mathrm{F})\) . Show that the following properties are equivalent:
  \end{enumerate}
\end{enumerate}

\begin{enumerate}
\def\labelenumi{(\roman{enumi})}
\item
  The image under \(u\) of every weakly compact subset of E is (strongly) compact;
\item
  The image under \(u\) of every weakly relatively compact subset of E is relatively compact in F;
\item
  The image under \(u\) of every weakly convergent sequence in E converges strongly in F.
\end{enumerate}

(Use the Šmulian theorem, EVT, IV, p. 36, th. 2.)

One says that u is semi-compact \(^{(3)}\) if it satisfies these properties.

\begin{enumerate}
\def\labelenumi{\alph{enumi})}
\setcounter{enumi}{1}
\item
  Let E and F be Banach spaces. Show that the set \(\mathcal{L}^{\mathrm{sc}}(\mathrm{E};\mathrm{F})\) of semi-compact mappings is a closed subspace of \(\mathcal{L}(\mathrm{E};\mathrm{F})\) .
\item
  Let E and F be Banach spaces and \(u \in \mathcal{L}(\mathrm{E};\mathrm{F})\) . Let \(E_{1}\) , \(F_{1}\) be Banach spaces, \(f \in \mathcal{L}(\mathrm{E}_{1};\mathrm{E})\) and \(g \in \mathcal{L}(\mathrm{F};\mathrm{F}_{1})\) . If \(u: E \to F\) is semi-compact, so is \(g \circ u \circ f\) .
\item
  Let E and F be Banach spaces and \(u \in \mathcal{L}(\mathrm{E};\mathrm{F})\) . If u is compact, it is semi-compact. Conversely, if E is reflexive and u is semi-compact, then u is compact.
\item
  Let E and F be Banach spaces. If the dual space E′ of E is of countable type, then every completely continuous linear map E → F is compact.
\item
  Let I be an infinite set; the identity map of \(\ell^{1}(\mathbf{I})\) is semi-compact, but is not weakly compact (cf. EVT, IV, p. 54, exerc. 15). g) The injection of \(\ell_{\mathrm{K}}^{1}(\mathbf{N})\) into \(\ell_{\mathrm{K}}^{2}(\mathbf{N})\) is semi-compact but is not compact (cf. loc. cit.).
\item
  Let E be a Banach space. Every continuous linear map \(E \rightarrow \ell_{K}^{1}(\mathbf{N})\) is semi-compact, and every continuous linear map \(\ell_{\mathrm{K}}^{1}(\mathbf{N}) \rightarrow \mathrm{E}\) is semi-compact.
(3) Some authors call them completely continuous maps.
\item
  Let E be a Banach space. If E is reflexive or if its dual is of countable type, then every continuous linear map \(\mathrm{E} \to \ell_{\mathrm{K}}^{1}(\mathbf{N})\) is compact.
\end{enumerate}

\begin{enumerate}
\def\labelenumi{\arabic{enumi})}
\setcounter{enumi}{30}
\tightlist
\item
  Let E be a Banach space.
\end{enumerate}

\begin{enumerate}
\def\labelenumi{\alph{enumi})}
\tightlist
\item
  Prove that the following conditions are equivalent:
\end{enumerate}

\begin{enumerate}
\def\labelenumi{(\roman{enumi})}
\item
  For any Banach space F, every weakly compact mapping from E into F is semi-compact.
\item
  If \((x_{n})\) is a sequence tending to 0 in E for the topology \(\sigma(\mathrm{E},\mathrm{E}^{\prime})\) and \((x_{n}^{\prime})\) a sequence tending to 0 in \(E^{\prime}\) for the topology \(\sigma(\mathrm{E}^{\prime},\mathrm{E}^{\prime\prime})\) , the sequence \(\langle x_{n},x_{n}^{\prime}\rangle\) tends to 0.
\end{enumerate}

(To prove that (i) implies (ii), one takes \(F = c_0\) . If (i) is not satisfied, there exists a weakly compact map \(u\) from E into a Banach space F and a sequence \((x_n)\) in E tending weakly to 0, such that \(\inf(\|u(x_n)\|) > 0\) ; one constructs a sequence \((y_n)\) in the unit ball of \(E'\) , convergent for \(\sigma(E', E'')\) and such that \(\langle u(x_n), y_n \rangle = \|u(x_n)\|\) , and one takes \(x_n' = y_n - \lim y_n\) .)

These conditions are satisfied in particular when \(E = \mathcal{C}(T)\) , where T is a compact space, and when \(E = L^{1}(X, \mu)\) , or \(E = L^{\infty}(X, \mu)\) , where X is a locally compact space and \(\mu\) a measure on X (INT, VI, § 2, exerc. 23). b) Prove that a Banach space possessing the above properties admits no reflexive direct factor of infinite dimension.

\begin{enumerate}
\def\labelenumi{\alph{enumi})}
\setcounter{enumi}{2}
\tightlist
\item
  Let H be a reflexive Banach space, and let B be the unit ball of H', equipped with the weak topology. Show that H can be identified with a closed non-direct subspace of \(\mathcal{C}(\mathrm{B})\) .
\end{enumerate}

\begin{enumerate}
\def\labelenumi{\arabic{enumi})}
\setcounter{enumi}{31}
\item
  A Banach space of countable type admitting a Schauder basis (EVT, IV, p. 70, exerc. 14) possesses the approximation property.
\item
  \begin{enumerate}
  \def\labelenumii{\alph{enumii})}
  \tightlist
  \item
    Prove that the strict inductive limit of a sequence of approximation spaces is an approximation space.
  \end{enumerate}
\end{enumerate}

\begin{enumerate}
\def\labelenumi{\alph{enumi})}
\setcounter{enumi}{1}
\tightlist
\item
  Prove that the locally convex direct sum of a family of approximation spaces is an approximation space.
\end{enumerate}

\begin{enumerate}
\def\labelenumi{\arabic{enumi})}
\setcounter{enumi}{33}
\tightlist
\item
  Let E be a Banach space. Prove that the following properties are equivalent:
\end{enumerate}

\begin{enumerate}
\def\labelenumi{(\roman{enumi})}
\item
  The space E' is an approximation space;
\item
  For every Banach space F, the closure of \(\mathcal{L}^{\mathrm{f}}(\mathrm{E};\mathrm{F})\) in \(\mathcal{L}(\mathrm{E};\mathrm{F})\) is \(\mathcal{L}^{\mathrm{c}}(\mathrm{E};\mathrm{F})\) ;
\item
  For every closed subspace F of \(c_{0}\) the closure of \(\mathcal{L}^{\mathrm{f}}(\mathrm{E};\mathrm{F})\) in \(\mathcal{L}(\mathrm{E};\mathrm{F})\) is \(\mathcal{L}^{\mathrm{c}}(\mathrm{E};\mathrm{F})\) .
\end{enumerate}

\begin{enumerate}
\def\labelenumi{\arabic{enumi})}
\setcounter{enumi}{34}
\tightlist
\item
  Let E be a Banach space. We say that E has the metric approximation property if the identity map of E belongs to the closure of the set of elements of norm \(\leqslant 1\) of \(\mathcal{L}^{\mathrm{f}}(\mathrm{E})\) in the space \(\mathcal{L}(\mathrm{E})\) endowed with the topology of compact convergence.
\end{enumerate}

Prove that the following properties are equivalent:

\begin{enumerate}
\def\labelenumi{(\roman{enumi})}
\item
  The space E has the metric approximation property;
\item
  The unit ball of \(\mathcal{L}(\mathrm{E})\) is the closure, for the topology of compact convergence, of the set of elements of norm \(\leqslant 1\) of \(\mathcal{L}^{\mathrm{f}}(\mathrm{E})\) ;
\item
  For every Banach space F, the unit ball of \(\mathcal{L}(\mathrm{E};\mathrm{F})\) is the closure, for the topology of compact convergence, of the set of elements of norm \(\leqslant 1\) of \(\mathcal{L}^{\mathrm{f}}(\mathrm{E};\mathrm{F})\) ;
\item
  For every Banach space F, the unit ball of \(\mathcal{L}(\mathrm{F};\mathrm{E})\) is the closure, for the topology of compact convergence, of the set of elements of norm \(\leqslant 1\) of \(\mathcal{L}^{\mathrm{f}}(\mathrm{F};\mathrm{E})\) .
\end{enumerate}

\begin{enumerate}
\def\labelenumi{\arabic{enumi})}
\setcounter{enumi}{35}
\tightlist
\item
  \begin{enumerate}
  \def\labelenumii{\alph{enumii})}
  \tightlist
  \item
    Let X be a locally compact topological space. Prove that the space \(\mathcal{C}_{0}(\mathrm{X})\) has the metric approximation property (adapt the proof of Prop. 13 of III, p. 21).
  \end{enumerate}
\end{enumerate}

\begin{enumerate}
\def\labelenumi{\alph{enumi})}
\setcounter{enumi}{1}
\tightlist
\item
  Let X be a locally compact topological space equipped with a positive measure and \(p \in [1, +\infty]\) . Prove that the space \(L^{p}(X)\) has the metric approximation property (same method).
\end{enumerate}

\begin{enumerate}
\def\labelenumi{\arabic{enumi})}
\setcounter{enumi}{36}
\tightlist
\item
  Let X be a differentiable manifold of class \(C^{r}\) where \(1 \leqslant r \leqslant \infty\) , paracompact, locally of finite dimension, and let E be a vector bundle over X of class \(C^{r}\) , locally of finite dimension. Prove that the space \(\mathcal{S}^{r}(\mathrm{X};\mathrm{E})\) has the approximation property.
\end{enumerate}

